%%%%%%%%%%%%%%%%%%%%%%%%%%%%%%%%%%%%%%%%%%%%%%%%%%%%%%%%%%%%%%%%%%%%%%
% Lớp 8 | Chương III. Tứ giác | Tứ giác
% Dạng 8H3?10-1: Nhận biết tứ giác, các yếu tố của tứ giác
%%%%%%%%%%%%%%%%%%%%%%%%%%%%%%%%%%%%%%%%%%%%%%%%%%%%%%%%%%%%%%%%%%%%%%


% nguồn: De_AI_Toan8_cuoi_hoc_ki_i_1 (nhập 21/11/2026)
\begin{bt}%[8H3N10-1]%[ma:43df6ed4]
	Cho tứ giác lồi $ABCD$. Hãy kể tên:
a) Các đỉnh, các cạnh và hai đường chéo của tứ giác;
b) Các cặp cạnh đối và các cặp góc đối của tứ giác.
	\loigiai{
		a) Tứ giác $ABCD$ có: \\
- Bốn đỉnh: $A$, $B$, $C$, $D$. \\
- Bốn cạnh: $AB$, $BC$, $CD$, $DA$. \\
- Hai đường chéo: $AC$ và $BD$. \\
b) Các cặp cạnh đối là: $AB$ và $CD$; $BC$ và $DA$. \\
Các cặp góc đối là: $\widehat{A}$ và $\widehat{C}$; $\widehat{B}$ và $\widehat{D}$.
	}
\end{bt}

% nguồn: De_AI_Toan8_cuoi_hoc_ki_i_1 (nhập 21/11/2026)
\begin{bt}%[8H3N10-1]%[ma:ce64cfa5]
	Hãy giải thích xem có tồn tại tứ giác lồi nào có số đo bốn góc thỏa mãn từng trường hợp sau hay không:
a) Bốn góc đều bằng $90^\circ$;
b) Bốn góc lần lượt bằng $100^\circ$, $110^\circ$, $80^\circ$, $80^\circ$;
c) Cả bốn góc đều là góc nhọn.
	\loigiai{
		a) Có tồn tại (ví dụ hình chữ nhật, hình vuông), vì $90^\circ + 90^\circ + 90^\circ + 90^\circ = 360^\circ$ và mỗi góc đều nhỏ hơn $180^\circ$. \\
b) Không tồn tại, vì tổng bốn góc là $100^\circ + 110^\circ + 80^\circ + 80^\circ = 370^\circ \ne 360^\circ$. \\
c) Không tồn tại, vì nếu cả bốn góc đều nhọn (nhỏ hơn $90^\circ$) thì tổng bốn góc sẽ nhỏ hơn $4 \cdot 90^\circ = 360^\circ$, trái với định lí tổng các góc của tứ giác.
	}
\end{bt}
